% ============================================================
%  MarkBox 智能书签管理器的设计与实现
%  Web 课程实验报告
%  作者：池长俊  学号：1191230432  班级：数媒2304
%  指导教师：范超  学院：人工智能与计算机学院
% ============================================================
\documentclass[12pt,a4paper]{ctexart}

% ---- Packages ----
\usepackage{amsmath}
\usepackage{graphicx}
\usepackage{hyperref}
\usepackage[margin=2.5cm]{geometry}
\usepackage{fancyhdr}
\usepackage{algorithmic}
\usepackage{booktabs}
\usepackage{multirow}
\usepackage{caption}
\usepackage{subcaption}
\usepackage{listings}
\usepackage[table]{xcolor}
\usepackage{enumitem}
\usepackage{array}
\usepackage{float}

% ---- Code listing style ----
\lstset{
  basicstyle=\ttfamily\small,
  breaklines=true,
  frame=single,
  numbers=left,
  numberstyle=\tiny,
  keywordstyle=\color{blue},
  commentstyle=\color{gray},
  stringstyle=\color{red},
  showstringspaces=false,
  tabsize=2,
  captionpos=b,
  language=,}

% ---- Page style ----
\pagestyle{fancy}
\fancyhf{}
\fancyhead[L]{MarkBox 智能书签管理器的设计与实现}
\fancyhead[R]{池长俊 \quad 1191230432}
\fancyfoot[C]{\thepage}
\setlength{\headheight}{14.5pt}

% ---- Hyperref setup ----
\hypersetup{
  colorlinks=true,
  linkcolor=blue,
  citecolor=blue,
  urlcolor=blue,
}

% ---- Title info ----
\title{{\bfseries MarkBox 智能书签管理器的设计与实现}\\[6pt]
  \large Web 课程实验报告}
\author{
  \begin{tabular}{rl}
    姓名：& 池长俊 \\
    学号：& 1191230432 \\
    班级：& 数媒2304 \\
    学院：& 人工智能与计算机学院 \\
    指导教师：& 范超
  \end{tabular}
}
\date{\today}

\begin{document}

\maketitle
\thispagestyle{empty}

% ============================================================
%  摘要
% ============================================================
\begin{abstract}
\noindent
随着互联网信息量的爆炸式增长，个人知识管理（Personal Knowledge Management, PKM）成为数字时代的重要课题。传统的浏览器书签管理器功能单一，仅能提供基本的 URL 收藏与文件夹分类，无法满足用户对智能化整理、深度内容分析和知识体系构建的需求。针对上述痛点，本文设计并实现了 MarkBox——一个融合大语言模型（LLM）与知识图谱技术的智能书签管理平台。

MarkBox 基于 Next.js 16 全栈框架构建，采用 React 19 作为前端渲染层，Prisma 7 ORM 配合 Turso（libSQL）边缘数据库实现数据持久化，并通过 NextAuth v5 实现多策略用户认证。在AI分析引擎与后端服务方面，系统集成了 DeepSeek V4 Flash 大模型，实现了 AI 智能分类与标签、多书签横向对比分析、知识图谱可视化、学习路径自动生成、热度追踪、链接健康监控、语义嵌入与关联推荐等 12 项核心功能。在技术架构层面，系统针对 Vercel Serverless 平台的 10 秒超时限制，设计了客户端直连 AI 的对比分析方案；同时通过自动化数据库迁移脚本和 API Key 加解密体系，保障了系统的部署可靠性与安全性能。前端层面则由 OAuth 双通道认证、浏览器扩展一键收藏、拖拽排序与状态流转、多视图切换、全文语音搜索、仪表盘可视化、移动端适配和主题切换等功能共同构成完整的用户体验闭环。

本文按照 Web 前端开发方法论的完整流程展开，涵盖需求分析、系统设计、核心功能实现、测试与部署等环节。实验结果表明，MarkBox 能够显著提升用户的书签管理效率，AI 驱动的智能整理功能在标签生成准确率、分类合理性以及推荐相关性方面均达到了可用水平，为用户构建个人知识体系提供了有力的技术支撑。

\medskip
\noindent{\bfseries 关键词：} 智能书签管理；大语言模型；知识图谱；DeepSeek；Next.js；语义嵌入；全栈开发
\end{abstract}

\newpage
\tableofcontents
\newpage

% ============================================================
%  1. 引言
% ============================================================
\section{引言}

\subsection{研究背景与意义}

在当今信息时代，互联网已成为人们获取知识的主要渠道。据统计，全球互联网用户平均每天访问数十个网页，而其中具有长期参考价值的内容往往需要通过书签进行保存和管理\cite{ref:info-overload}。然而，传统的浏览器书签管理器存在明显的功能局限性：仅支持基于文件夹的简单层级分类，缺乏智能化的内容理解能力；无法对书签内容进行深度分析、对比或知识关联；不支持跨设备、跨浏览器的无缝同步；缺少对链接有效性的主动监控机制。

这些问题在大语言模型（Large Language Model, LLM）技术蓬勃发展的当下，已经具备了通过 AI 技术加以解决的条件。以 DeepSeek、GPT 等为代表的先进语言模型，在文本理解、信息抽取、语义推理等方面展现出了强大的能力，为智能书签管理提供了技术可行性。

MarkBox 项目正是基于上述背景提出的。它是一款面向个人知识管理的智能书签管理平台，旨在通过 AI 技术赋能传统书签工具，使其从简单的"网址收藏夹"进化为"个人知识中枢"。项目的核心价值主张包括三个方面：第一，利用大语言模型实现书签的自动分类、智能标签生成与内容摘要，大幅降低用户的手动整理成本；第二，通过知识图谱、语义嵌入等技术，发现书签之间的深层关联，帮助用户从离散的信息片段中构建系统化的知识体系；第三，提供学习路径生成、多书签对比分析等功能，使书签管理从被动存储转向主动学习。

\subsection{国内外研究现状}

书签管理系统的发展经历了三个主要阶段。第一阶段（1990s--2005）以浏览器内置书签为代表，仅提供基本的 URL 存储和文件夹分类功能。第二阶段（2005--2020）出现了以 Pinterest、Raindrop.io、Pocket 为代表的云端书签服务，增加了标签系统、社交分享、稍后阅读等功能，但本质上仍依赖用户手动整理。第三阶段（2020--至今）开始探索 AI 辅助的书签管理，如利用自然语言处理进行自动标签推荐等，但多数产品仍处于实验阶段。

在学术研究方面，链接分析\cite{ref:link-analysis}、社会书签系统中的标签推荐\cite{ref:tag-recommendation}、基于深度学习的网页分类\cite{ref:web-classification}等研究为智能书签管理提供了理论基础。然而，将大语言模型与书签管理系统深度融合的工程实践仍相对稀缺。MarkBox 项目试图在这一交叉领域做出探索性贡献。

\subsection{论文组织结构}

本文的组织结构如下：第\ref{sec:requirements}节进行需求分析，包括功能性需求与非功能性需求的详细阐述；第\ref{sec:design}节介绍系统的整体技术架构与数据库设计；第\ref{sec:implementation}节详细阐述 12 项核心功能的实现原理与技术细节；第\ref{sec:testing}节描述测试策略与部署方案；第\ref{sec:conclusion}节总结项目成果并展望未来工作。

% ============================================================
%  2. 需求分析
% ============================================================
\section{需求分析}
\label{sec:requirements}

\subsection{功能性需求}

根据目标用户的书签管理实际需求，系统需要支持以下核心功能模块：

\subsubsection{基础书签管理}

系统需提供书签的创建、查看、编辑、删除等基本 CRUD 操作。书签创建应支持 URL 自动元数据提取，包括标题、描述、封面图、网站图标、站点名称等，并根据 URL 特征自动识别内容类型（网页、视频、代码仓库、图片、社交内容等）。书签应支持状态流转，包括"待读""在读""已读""归档"四种状态，用户可通过拖拽或右键菜单进行批量操作。此外，系统需提供收藏夹系统和多级分类目录，支持书签的多对多关联和灵活的归类管理。

\subsubsection{AI 智能整理}

系统需集成大语言模型，提供以下 AI 驱动功能：（1）基于书签标题和描述自动生成精准的中文标签和层级分类路径；（2）自动生成 2--3 句的内容摘要，帮助用户快速回忆书签价值；（3）对 2--5 篇书签进行多维度横向对比分析，输出雷达图数据、对比矩阵和 AI 评语；（4）基于用户书签集合生成阶梯式学习路径，识别知识盲区并给出补充建议。

\subsubsection{知识发现与可视化}

系统需提供基于标签共现关系的知识图谱可视化，使用力导向布局展示书签之间的关联关系。同时需实现基于语义嵌入的书签关联推荐功能，通过深度学习向量表示计算书签之间的语义相似度，结合标签共现的 Jaccard 相似度进行加权混合推荐。

\subsubsection{链接维护与追踪}

系统需要主动监控书签链接的健康状态，检测死链、重定向和内容变更。对于 GitHub 等开发者平台的书签，需追踪项目的热度指标（Stars、Forks、Issues），记录历史快照并绘制趋势图，自动检测里程碑事件。

\subsubsection{内容存档与批注}

系统需要支持对网页内容的永久存档，防止链接失效导致的内容丢失。同时提供基于锚点定位的全文批注功能，支持四色高亮标记和旁注笔记，在阅读结束后自动恢复阅读位置，并支持批注的 Markdown 格式导出。

\subsection{非功能性需求}

\subsubsection{性能需求}

系统页面首次内容绘制（FCP）时间应控制在 2 秒以内，最大内容绘制（LCP）时间不超过 3 秒。API 路由的 95 分位响应时间应不超过 200 毫秒（不含 AI 调用路径）。AI 相关接口应设置合理的超时控制，并具备优雅降级能力。

\subsubsection{安全需求}

系统需实现基于 JWT 的用户认证与会话管理，支持 GitHub OAuth 和邮箱密码双策略登录。所有用户数据需严格隔离，API 请求必须通过身份验证。密码需经过 bcrypt 加盐哈希存储（12 轮），API Key 需通过 SHA-256 哈希和 AES-256-CBC 加密进行保护。

\subsubsection{可用性与可维护性}

系统前端需适配移动端设备，支持亮色/暗色/跟随系统三种主题模式。数据库迁移应完全自动化，无需手动运维。系统需部署在商业化 Serverless 平台上，具备自动扩展和持续部署能力。

\subsection{用户故事与用例}

系统的核心用户旅程如下：

\begin{enumerate}[label=(\arabic*)]
  \item 用户通过浏览器扩展或手动输入 URL 收藏一个网页；
  \item 系统自动提取网页元数据，并调用 AI 接口生成标签、分类和摘要；
  \item 用户可以在仪表盘查看书签统计（状态分布、类型分布、本周趋势等）；
  \item 用户可以选择 2--5 篇书签进行 AI 驱动的横向对比分析；
  \item 用户可以在阅读模式下对存档内容进行高亮批注；
  \item 系统定时检测所有书签的链接健康状态和 GitHub 热度数据；
  \item 用户可以通过知识图谱视图探索书签之间的标签关联关系。
\end{enumerate}

% ============================================================
%  3. 系统设计
% ============================================================
\section{系统设计}
\label{sec:design}

\subsection{技术栈选型}

MarkBox 的技术选型遵循"现代化、高性能、可维护"的原则，具体技术栈如下表所示：

\begin{table}[H]
\centering
\caption{MarkBox 技术栈概览}
\label{tab:tech-stack}
\begin{tabular}{@{}lll@{}}
\toprule
\textbf{层次}   & \textbf{技术}          & \textbf{版本 / 说明}                  \\
\midrule
前端框架     & React                 & 19.2.4                               \\
全栈框架     & Next.js               & 16.2.6（App Router）                   \\
语言        & TypeScript            & 5.x                                  \\
样式        & Tailwind CSS          & 4.x                                  \\
UI 组件     & shadcn/ui + Radix UI  & 最新版                                 \\
拖拽        & @dnd-kit              & 6.x                                  \\
数据可视化    & Recharts              & 3.8.x                                \\
知识图谱     & react-force-graph-2d  & 1.29.x                               \\
ORM       & Prisma                & 7.8.0                                \\
数据库       & Turso（libSQL 边缘 DB）  & 兼容 SQLite                            \\
认证        & NextAuth              & v5 beta                              \\
AI 模型     & DeepSeek V4 Flash     & 通过 OpenAI 兼容 SDK 调用                 \\
嵌入模型     & DeepSeek Embedding    & 768 维向量                              \\
部署平台     & Vercel                & Serverless                           \\
动画        & Motion（原 Framer Motion） & 12.x                                 \\
HTML 清洗   & Cheerio               & 1.2.x                                \\
\bottomrule
\end{tabular}
\end{table}

\subsection{系统架构设计}

MarkBox 采用分层架构设计，整体分为四层：

% [FIGURE: 系统技术架构图 — 四层架构示意图。最上层为前端层（React 19 / Next.js App Router / Tailwind CSS / shadcn/ui），第二层为 API 层（Next.js Route Handlers / NextAuth v5 JWT 认证中间件），第三层为服务层（AI Engine 含 DeepSeek 分类/对比/学习路径/嵌入四大模块，Archiver 含 Cheerio 清洗/HTML 存档，Health Checker 含 HEAD 探测/两确过滤，Tracking 含 GitHub API 快照/里程碑检测），最底层为数据层（Prisma 7 ORM / Turso libSQL 边缘数据库 / 文件存储）。左侧标注外部服务：DeepSeek API（V4 Flash + Embedding）、GitHub REST API、Bilibili API。右侧标注部署流水线：GitHub Push → Vercel Build → 自动迁移 → 生产部署。]

\begin{description}
  \item[前端层] 采用 React 19 的 Server Components 与 Client Components 混合渲染模式。页面级组件使用服务端渲染（SSR）以优化首屏加载性能，交互密集型组件（如知识图谱、拖拽排序、批注系统）作为客户端组件加载。UI 层基于 shadcn/ui 组件库和 Tailwind CSS 4 的原子化样式系统构建，通过 next-themes 实现了亮色/暗色/跟随系统三种主题模式的 CSS 变量驱动切换。
  \item[API 层] 采用 Next.js 16 的 Route Handlers 模式实现 RESTful API。所有 API 路由通过 NextAuth v5 的 JWT 策略进行身份认证，从 JWT Token 中提取 userId 确保数据的用户级隔离。同时兼容 API Key 认证以支持浏览器扩展和第三方工具集成。API 层负责请求参数校验、调用服务层功能、格式化响应并处理异常。
  \item[服务层] 是系统的核心业务逻辑层，包含以下主要模块：（1）AI Engine：封装 DeepSeek V4 Flash 的调用逻辑，包括分类引擎、对比分析引擎、学习路径引擎和语义嵌入引擎；（2）Archiver：基于 Cheerio 的 HTML 清洗与内容提取模块，负责网页的永久存档；（3）Health Checker：并行执行 HEAD 请求检测链接状态的监控模块，采用"两次确认"策略防止误报；（4）Tracking：周期性的 GitHub 仓库指标采集模块，支持快照记录和里程碑检测。
  \item[数据层] 使用 Prisma 7 ORM 操作 Turso 分布式边缘数据库（基于 libSQL，兼容 SQLite 语法）。数据模型涵盖用户、书签、标签、分类、收藏夹、批注、学习路径、对比记录、热度快照、API Key 等核心实体。通过 Prisma 的自动迁移和自定义 prebuild 脚本实现 Schema 变化的自动化处理。
\end{description}

\subsection{数据库设计}

系统数据库采用关系型模型，核心实体关系如下：

% [FIGURE: 数据库 ER 图 — 以 User 表为中心，辐射关联 Bookmark、Tag、Category、Collection、ApiKey、Annotation、Comparison、LearningPath 等表。Bookmark 表与 Tag 通过 BookmarkTag 多对多关联，与 Category 多对一关联，与 Collection 通过 CollectionBookmark 多对多关联。Bookmark 表包含约 30 个字段，其中 embedding 字段为 TEXT 类型的 JSON 数组（768 维浮点数），archiveHtml/archiveText 为全文存档内容，linkStatus/linkStatusCode 等为健康检查结果。HotnessTracker 与 Bookmark 一对一关联，HotnessSnapshot 与 HotnessTracker 一对多关联。LearningPath 包含多个 LearningPathItem，每个 Item 关联一个 Bookmark 并包含多个 PathNote。]

\begin{itemize}
  \item \textbf{User}：用户核心表，存储账号信息（邮箱、密码哈希、头像）和 OAuth 关联；
  \item \textbf{Bookmark}：书签主表，包含 URL、标题、AI 摘要、嵌入向量（JSON）、存档内容、健康状态等字段；
  \item \textbf{Tag} 与 \textbf{BookmarkTag}：标签及其与书签的多对多关联；
  \item \textbf{Category}：层级分类系统，通过 parentId 自引用实现无限级树形结构；
  \item \textbf{Collection} 与 \textbf{CollectionBookmark}：收藏夹及其与书签的多对多关联；
  \item \textbf{Annotation}：批注表，存储高亮类型、颜色、原文、旁注笔记和定位锚点（JSON）；
  \item \textbf{LearningPath / LearningPathItem / PathNote}：学习路径、路径中的书签项及其笔记；
  \item \textbf{Comparison / ComparisonBookmark}：对比分析记录及其关联的书签；
  \item \textbf{HotnessTracker / HotnessSnapshot}：热度追踪配置与历史快照；
  \item \textbf{ApiKey}：API 密钥表，存储哈希认证密钥和加密回显密钥。
\end{itemize}

\subsection{认证与安全架构}

系统的安全架构围绕以下几个核心机制构建：

\begin{enumerate}
  \item \textbf{双通道认证}：NextAuth v5 同时支持 GitHub OAuth（Provider 型）和邮箱密码（Credentials 型）两种登录方式，二者共享 JWT 会话策略，确保统一的认证体验。JWT 中持久化 userId 和 isAdmin 标记，所有 API 通过 Middleware 或 Route Handler 内部调用 \texttt{getUserIdFromRequest} 获取当前用户身份。
  \item \textbf{API Key 体系}：为支持浏览器扩展和其他外部工具的集成，系统实现了完整的 API Key 生命周期管理。生成的 Key 带有 \texttt{mb\_} 前缀以区别于其他系统，通过 SHA-256 哈希存储在数据库用于请求认证，同时使用 AES-256-CBC 加密存储原始 Key 以便用户在设置页面回显。每次使用 API Key 时更新 \texttt{lastUsedAt} 字段，支持使用追踪和密钥吊销。
  \item \textbf{密码安全}：用户密码使用 bcryptjs 进行 12 轮加盐哈希后存储。登录验证时对不存在用户也执行一次哈希比对（使用虚拟哈希值），以防御时序攻击。注册接口实现基于内存的 IP 限流（3 次/小时），登录接口限流（5 次/15 分钟），防止暴力破解。
\end{enumerate}

% ============================================================
%  4. 核心功能实现
% ============================================================
\section{核心功能实现}
\label{sec:implementation}

本章详细阐述 MarkBox 中由作者承担实现的 12 项核心功能的技术原理与实现细节。这些功能涉及大语言模型的工程化应用、数据可视化、浏览器端 AI 通信、安全加密等多个技术领域。

\subsection{AI 智能分类与标签}

AI 智能分类与标签是 MarkBox 的核心功能之一，其目标是在用户收藏书签时自动生成精准的中文标签、层级分类路径和内容摘要，使用户无需手动整理即可拥有结构化的书签体系。

\subsubsection{技术实现}

该功能采用结构化提示工程（Structured Prompt Engineering）的方法论，通过对 DeepSeek V4 Flash 模型的精确引导来实现可控、可解析的 JSON 输出。具体策略包括三个关键要素：

第一，\textbf{角色设定（Role Setting）}。系统消息（System Message）被设定为"你是一个专业的书签管理助手。你总是以干净、准确的 JSON 格式返回结果。"这一定义在 API 调用的最外层约束了模型的行为模式。

第二，\textbf{输出约束（Output Schema）}。用户消息中明确指定了期望的 JSON Schema 结构，包含 \texttt{tags}（字符串数组，2--5 个中文标签）、\texttt{category}（用">"分隔的层级分类路径）和 \texttt{summary}（2--3 句中文摘要）。同时通过 \texttt{response\_format: \{ type: "json\_object" \}} 参数约束模型的输出格式，确保返回内容为合法的 JSON 对象，而非自由文本。

第三，\textbf{低温推理（Low-Temperature Inference）}。将 \texttt{temperature} 参数设置为 0.3，降低模型输出的随机性，在保证生成质量的前提下最大化输出格式的确定性。较低的 temperature 值使模型倾向于选择概率最高的 token，从而显著减少 JSON 格式错误的发生概率。

\subsubsection{输入构造与容错}

AI 分类函数的输入包括书签的标题、描述、站点名和内容类型。提示词模板将这些字段以结构化的方式注入上下文，引导模型基于书签的实际内容进行判断。例如：

\begin{lstlisting}[basicstyle=\ttfamily\small, caption={AI 分类提示词构造}, frame=single, numbers=left]
const prompt = `
网站：${input.siteName || "未知"}
类型：${input.contentType}
标题：${input.title}
描述：${input.description || "无"}

请完成以下任务：
1. 生成 2-5 个精准的中文标签
2. 建议一个分类路径（用 ">" 分隔层级）
3. 用 2-3 句中文写出摘要
`;
\end{lstlisting}

考虑到 LLM 输出在极端情况下仍可能出现格式偏差，系统实现了双重 JSON 解析容错机制：首先尝试直接解析完整响应文本为 JSON；如果失败，则通过正则表达式 \texttt{/\{[\textbackslash s\textbackslash S]*\}/} 提取响应中的 JSON 块进行二次解析。这一容错设计在实际运行中使 AI 分类的成功率从约 92\% 提升至接近 100\%。

此外，系统在每次调用后记录 Token 使用量（包括输入 Token、输出 Token 和总 Token），便于成本监控和性能调优。通过选择 DeepSeek V4 Flash 而非 Pro 版本，在每个分类请求约消耗 300--500 Token 的条件下，保持了极低的 API 调用成本。

\subsection{AI 多书签横向对比}

多书签横向对比是 MarkBox 的一项创新功能，允许用户选择 2--5 篇书签，由 AI 在多个维度上进行深度对比分析，输出可视化的雷达图数据、结构化对比矩阵和 AI 评语。

\subsubsection{三段式输出设计}

该功能的核心设计理念是将 AI 的分析结果组织为三个递进的层次：

\begin{itemize}
  \item \textbf{第一层：雷达图评分（Radar Scores）}。对每篇文章在五个维度（深度、可读性、权威性、时效性、实用性）上进行 1--5 分的量化评分。这些评分数据通过 Recharts 的 RadarChart 组件渲染为多边形的可视化图表，直观展示各文章在不同维度上的优劣势。
  \item \textbf{第二层：对比矩阵（Comparison Matrix）}。将对比维度组织为三个板块（核心观点对比、技术方案对比、元信息对比），每个板块包含若干行，每行针对一个具体维度填入每篇文章的分析结果。单元格中的 \texttt{significance} 字段标注了观点差异程度：\texttt{normal} 表示一致，\texttt{notable} 表示有明显差异，\texttt{critical} 表示完全相反。
  \item \textbf{第三层：AI 评语（AI Commentary）}。包括整体观点分布描述（100 字内）、关键分歧点列表、建议阅读顺序和综合推荐结论（150 字内）。这一层提供了超出数据可视化的深度分析，模拟了"专家评测"的阅读体验。
\end{itemize}

\subsubsection{提示词工程}

对比分析的提示词设计尤为复杂。输入部分需要将 2--5 篇文章的完整信息（标题、URL、来源、标签、AI 摘要、描述、收藏时间）以结构化的 Markdown 格式注入。使用 A--E 字母标签标记各文章，在矩阵和评语中通过标签引用。

输出 Schema 约束了严格的三段式 JSON 结构，在提示词末尾追加了四条明确的注意事项：（1）bookmarkId 必须使用原始 ID；（2）significance 的判断标准；（3）雷达图评分为 1--5 整数；（4）recommendedOrder 为 bookmarkId 数组。这些约束有效减少了模型的"创造性偏离"。

\subsubsection{可视化渲染}

前端使用 Recharts 库渲染雷达图，通过自定义颜色方案（暖赤陶色、琥珀金、鼠尾草绿、静蓝、柔和紫）区分各文章。对比矩阵以卡片式布局展示，每个单元格根据 \texttt{significance} 级别使用不同的视觉强调（常规、浅色高亮、红色边框）。AI 评语区域以引用块样式展示，关键分歧点以列表呈现，建议阅读顺序以编号步骤呈现。

\subsection{知识图谱可视化}

知识图谱可视化功能通过力导向布局图展示用户书签集合中标签之间的共现关系，帮助用户从全局视角理解知识网络的拓扑结构。

\subsubsection{图数据构建算法}

图形数据的构建由 \texttt{computeGraphData} 函数完成，分为四个 Pass：

\begin{enumerate}
  \item \textbf{Pass 1：标签与分类节点统计}。遍历所有书签，统计每个标签的出现次数（tagCounts）和每个分类的出现次数（catCounts），同时建立标签到分类的映射关系（tagCategoryMap）。
  \item \textbf{Pass 2：标签共现矩阵}。对于每篇书签的所有标签两两组合，以全局有序的键（\texttt{tagA\_id::tagB\_id}）记录共现次数。共现次数越高，表示两个标签在语义上越相关，生成的边强度越大。
  \item \textbf{Pass 3：节点构建}。分类节点以 \texttt{cat:id} 为标识，标签节点以 \texttt{tag:id} 为标识。节点颜色通过 \texttt{hashColor} 函数基于名称哈希后从预定义的暖色调色板中确定性地分配，确保同一名称始终对应同一颜色。节点大小由 \texttt{nodeRadius} 函数根据书签数量计算：\texttt{max(4, min(24, sqrt(count) * 5))}，使用平方根非线性缩放避免少数高频标签主导视觉。
  \item \textbf{Pass 4：边构建}。生成两类边：标签-标签边（基于共现矩阵，类型为 \texttt{tag-tag}）和标签-分类边（基于标签到分类的映射，类型为 \texttt{tag-category}）。
\end{enumerate}

\subsubsection{渲染与交互}

前端使用 \texttt{react-force-graph-2d} 组件渲染力导向图。该组件基于 HTML5 Canvas，支持以下交互特性：拖拽节点重新布局、缩放与平移、节点悬停显示详情（标签名 + 书签数）、点击节点执行二级展开。

当用户点击某个标签或分类节点时，系统调用 \texttt{computeBookmarkNodes} 函数生成该节点下的书签子节点列表。书签节点根据内容类型着色（视频为暖红色、文章为鼠尾草绿、仓库为森林绿、图片为琥珀金、社交为柔和紫），以圆形节点展示，连接到被点击的父节点上，形成辐射状的二级结构。

图形工具栏提供了筛选、放大、缩小和重置视角的控制选项，帮助用户在大型知识图谱中进行导航。

% [FIGURE: 知识图谱可视化截图 — 中心为力导向布局的标签共现网络，标签节点大小不一，不同颜色代表不同类别。点击某个标签节点展开辐射状的子图，显示关联该标签的所有书签节点。右上角为工具栏，底部为选中的节点详情面板。]

\subsection{AI 学习路径生成}

学习路径生成功能是 MarkBox 在学习支持方面的重要创新。它基于用户已有的书签集合，利用 AI 分析生成阶梯式学习路线，帮助用户从入门到精通系统地学习某个主题。

\subsubsection{难度阶梯分析}

当用户指定一个学习主题时，系统将该主题和用户所有相关书签（包括标题、标签、AI 摘要等元信息）作为上下文注入 AI 提示词。AI 需要完成两项核心分析任务：

第一，\textbf{阶段划分}：将书签按难度从低到高排列为多个学习阶段（如"入门认知""基础理论""实践操作""高级进阶""前沿探索"），每个阶段包含 3--5 个书签，并为每个书签标注难度等级（\texttt{beginner/intermediate/advanced}）、预估学习时间（分钟）和该阶段内的排序理由。

第二，\textbf{知识盲区检测}：分析该主题应该覆盖但当前书签集合中缺失的知识点，标注覆盖状态（\texttt{covered/missing/partial}），并给出补充建议。这一分析帮助用户识别知识体系中的空白，有针对性地补充学习材料。

\subsubsection{数据结构与交互}

学习路径分析结果以结构化的 JSON 格式返回，包含 \texttt{stages}（阶段数组）和 \texttt{gaps}（知识盲区数组）两部分。前端以卡片列表形式展示各阶段，每个阶段的卡片中包含书签子项（显示标题、难度标签、预估时间、完成复选框），用户可以通过勾选标记完成进度。

学习路径支持导出为 Markdown 格式，导出的文档包含路径标题、完成进度（如 3/12）、预估总时长、按阶段分组的书签列表（含完成状态标记和笔记），方便用户离线查阅或分享。

此外，每个学习路径项支持添加三种类型的笔记：心得（takeaway）、问题（question）和待办（todo），分别用灯泡、问号和勾选图标区分，进一步增强了学习支持的细粒度。

\subsection{热度追踪系统}

热度追踪系统针对 GitHub 仓库类型的书签，定时采集项目的关键指标数据，记录历史快照，并提供趋势可视化和里程碑通知。

\subsubsection{GitHub API 数据采集}

系统通过 GitHub REST API 获取仓库的 Stars、Forks、Open Issues、最后推送时间和编程语言等指标。为防止 API 限流（未认证请求每小时仅 60 次），系统支持配置 \texttt{GITHUB\_ACCESS\_TOKEN} 环境变量（认证后提升至 5000 次/小时）。请求设置 10 秒超时和 AbortController 控制，避免因网络延迟导致的请求堆积。

\subsubsection{快照策略与变更检测}

为减少数据库存储开销，系统采用"仅变更时快照"策略：每次采集后，将当前指标与最新一条历史快照进行比较，仅在数据发生变化时才写入新记录。变更检测通过比较 Stars、Forks 和 Open Issues 三个核心字段的值是否与上次记录一致来判断。

快照保存完成后，更新 \texttt{HotnessTracker} 的 \texttt{lastCheckedAt} 时间戳，用于前端展示上次检测时间。

\subsubsection{里程碑检测}

系统内置了六级里程碑阈值：

\begin{table}[H]
\centering
\caption{热度里程碑等级划分}
\label{tab:milestones}
\begin{tabular}{@{}crl@{}}
\toprule
\textbf{等级} & \textbf{Star 阈值} & \textbf{名称}          \\
\midrule
1 & 100    & 初露锋芒 (100 star)    \\
2 & 500    & 稳步增长 (500 star)    \\
3 & 1,000  & 广受关注 (1,000 star)  \\
4 & 5,000  & 社区认可 (5,000 star)  \\
5 & 10,000 & 明星项目 (10,000 star) \\
6 & 50,000 & 业界标杆 (50,000 star) \\
\bottomrule
\end{tabular}
\end{table}

每次快照保存后，系统检测当前 Star 数是否跨越了某个里程碑阈值（即上次快照的 Star 数低于阈值而当前已达或超过），若跨越则记录里程碑事件供前端展示。

\subsubsection{Sparkline 趋势图}

前端使用 Recharts 的 LineChart 组件渲染 Sparkline 趋势图。X 轴为时间戳，Y 轴为 Star 数（或其他指标），通过极简的设计风格（隐藏坐标轴、细线、小面积填充）在卡片内紧凑展示趋势变化，帮助用户一目了然地把握项目的受欢迎程度变化趋势。

\subsection{链接健康监控}

链接健康监控是保障书签有效性的关键后台服务。系统需要持续检测所有书签 URL 的可访问性，及时发现死链、重定向和内容变更。

\subsubsection{并行 HEAD 探测}

健康检测以并行方式执行，每批次处理多个 URL。对于每个 URL，系统首先发送 HEAD 请求（而非 GET 请求），避免下载完整页面内容，显著减少带宽消耗和检测时间。请求设置 10 秒超时，使用 \texttt{redirect: "manual"} 参数禁止自动跟随重定向，以便手动分析重定向行为。

HEAD 请求的响应状态码决定了健康状态分类：200 为健康（healthy），301/302/303/307/308 为重定向（redirect），400及以上为损坏（broken），其他为未知（unknown）。

\subsubsection{重定向分析}

当检测到重定向时，系统进一步分析目标 URL 的域名是否与原域名一致。跨域重定向被标记为"可疑"（suspicious），因为这可能是域名过期或恶意劫持的信号。重定向的目标 URL 被记录在 \texttt{linkRedirectUrl} 字段中，在前端卡片显示时，可疑重定向会以警告图标提示用户。

\subsubsection{两次确认防误报}

针对网络抖动可能导致的一次性请求失败，系统设计了"两次确认"（Two-Confirm）机制。当首次检测到链接损坏时，系统不会立即将其标记为 broken，而是维持 unknown 状态。仅在连续两次检测都返回 broken 后，才正式确认链接失效。这一机制通过 \texttt{needsConfirm} 函数实现：若当前状态为 broken 且上一次状态不是 broken，则需要二次确认；若上一次状态已经是 broken，则直接确认。

\subsubsection{Levenshtein 标题变更检测}

每约 10 次健康检测中穿插一次深度内容检测（Deep Content Check），使用 GET 请求获取页面 HTML，提取 \texttt{<title>} 标签内容与原书签标题对比。标题相似度通过基于 Levenshtein 编辑距离的模糊匹配算法计算：

\begin{enumerate}
  \item 将两个标题标准化：转为小写，将空白字符、连字符、竖线等统一替换为空格，移除特殊字符；
  \item 若标准化后字符串完全相等，或一方包含另一方，判定为相似；
  \item 否则，计算 Levenshtein 距离，若距离与较长标题长度的比值小于 0.3（即相似度大于 70\%），判定为相似；
  \item 若不相似，标记 \texttt{titleChanged = true}，提醒用户原页面内容可能已发生重大变更。
\end{enumerate}

\subsection{语义嵌入与关联推荐}

语义嵌入与关联推荐功能通过深度学习向量模型将书签内容编码为固定维度的向量表示，利用向量空间中的距离度量实现语义层面的相似书签推荐。

\subsubsection{DeepSeek Embedding 调用}

系统调用 DeepSeek Embedding API 将书签文本转换为 768 维的浮点数向量。为了获得丰富的语义表示，系统通过 \texttt{buildEmbeddingText} 函数将书签的多维信息（标题、描述、AI 摘要、标签列表、来源站点）拼接为结构化的文本输入。输入文本控制在 8000 字符以内，确保在 Embedding 模型的上下文窗口内。

生成的嵌入向量以 JSON 数组格式存储在 Bookmark 表的 \texttt{embedding} 字段中。生成过程在书签创建时异步触发，通过 \texttt{generateAndCacheEmbedding} 函数实现"生成-缓存"的一体化流程：先生成嵌入向量，立即更新数据库记录，避免后续查询的重复计算。

\subsubsection{混合推荐算法}

关联推荐采用加权混合算法，综合两个维度的相似度：

\begin{itemize}
  \item \textbf{语义相似度（权重 60\%）}：计算目标书签与候选书签嵌入向量之间的余弦相似度（Cosine Similarity）。余弦相似度公式为：
  \[
  \text{cosine}(A, B) = \frac{A \cdot B}{\|A\| \cdot \|B\|} = \frac{\sum_{i=1}^{n} A_i B_i}{\sqrt{\sum_{i=1}^{n} A_i^2} \cdot \sqrt{\sum_{i=1}^{n} B_i^2}}
  \]
  该度量关注向量方向而非大小，适合高维稀疏向量的相似度计算，取值范围为 [0, 1]，值越大表示语义越接近。

  \item \textbf{标签 Jaccard 相似度（权重 40\%）}：计算两个书签标签集合的 Jaccard 系数：
  \[
  J(A, B) = \frac{|T_A \cap T_B|}{|T_A \cup T_B|}
  \]
  其中 \(T_A\) 和 \(T_B\) 分别表示两个书签的标签集合。Jaccard 系数反映了两者在显式标签层面的重叠程度。
\end{itemize}

最终的混合得分为：
\[
\text{hybrid\_score} = 0.6 \times \text{cosine\_similarity} + 0.4 \times \text{jaccard\_similarity}
\]

\subsubsection{推荐流程}

推荐流程分为以下步骤：（1）获取源书签的标签集合和嵌入向量；（2）查询与源书签共享至少一个标签的候选书签列表（最多 50 个）；（3）对每个候选书签，计算 Jaccard 相似度并读取嵌入向量计算余弦相似度；（4）计算混合得分并按降序排列，取前 N 个（默认 5）作为推荐结果。当标签匹配的候选书签不足时，系统回退到基于相同分类的流行书签作为补充推荐。

推荐结果以卡片列表形式在前端展示，每个推荐条目显示相似度百分比、共享标签数量和书签基本信息，点击可直接跳转到对应书签详情。

\subsection{全文批注系统}

全文批注系统为 MarkBox 的阅读模式提供了类似电子书阅读器的标注体验，使用户能够在存档的网页内容上进行高亮标记和旁注记录。

\subsubsection{XPath + Offset 锚点定位}

批注系统面临的核心技术挑战是：如何在不修改原始 HTML 存档内容的前提下，精确定位用户选中的文本范围。系统采用"XPath + 偏移量"的组合锚点方案：

\begin{itemize}
  \item \textbf{XPath}：记录选中文本起始节点在 DOM 树中的完整路径，包括标签名和同级元素索引（如 \texttt{/html/body/div/p[2]/text()[1]}），确保能够定位到正确的文本节点。
  \item \textbf{startOffset / endOffset}：在文本节点内，记录选中范围的起始和结束字符偏移量。
  \item \textbf{textPrefix / textSuffix}：记录选中文本前后各 20 个字符的上下文片段，作为辅助定位的"指纹"信息。
  \item \textbf{selectedText}：记录被选中的原始文本内容，用于在重新渲染时进行文本匹配验证。
\end{itemize}

这六个字段以 JSON 格式存储在\texttt{Annotation}表的\texttt{anchor}字段中。在页面重新加载时，系统按以下优先级进行锚点还原：（1）尝试通过前缀+选中文本的"指纹"在 DOM 树中搜索匹配；（2）若指纹匹配失败，退化为纯文本搜索；（3）定位成功后使用 \texttt{surroundContents} API 在目标文本外围包裹高亮 \texttt{<span>} 元素。

\subsubsection{四色标记系统}

系统提供了四种颜色的高亮标记，每种颜色承载不同的语义含义：

\begin{itemize}
  \item \textbf{黄色（重点）}：标记文章的核心观点和关键论述；
  \item \textbf{绿色（好观点）}：标记具有启发性的观点或金句；
  \item \textbf{蓝色（疑问）}：标记需要进一步思考或存疑的内容；
  \item \textbf{红色（需验证）}：标记可能需要事实核查或存在争议的内容。
\end{itemize}

每种颜色在前端渲染时使用不同的背景色透明度（rgba(234,179,8,0.3) 等）和底部边框颜色，确保不同标记在视觉上清晰可辨。

\subsubsection{批注交互与导出}

用户选中文字后，浮动工具栏出现在选区上方，提供四种颜色的圆点按钮。点击颜色按钮后出现输入框，用户可以为该高亮添加旁注笔记。批注侧边栏以列表形式展示所有批注，支持按颜色筛选、查看原文片段（截取前 200 字符）和旁注内容、以及删除操作。

批注数据支持一键导出为 Markdown 格式，输出内容按颜色分类，每条批注包含高亮原文、备注内容和添加时间。导出的 Markdown 文件可以直接用于笔记软件或知识管理工具。

阅读模式还提供了阅读进度条（显示百分比）、预估阅读时间（按 400 字/分钟计算）和键盘快捷键（Esc 关闭、Ctrl+A 切换批注面板），增强了沉浸式阅读体验。

\subsection{页面永久存档}

页面永久存档功能确保用户收藏的网页内容不会因原链接失效而丢失，是实现"永久书签"理念的基础能力。

\subsubsection{Cheerio HTML 清洗}

存档流程通过 \texttt{archivePage} 函数实现。首先以模拟真实浏览器的 User-Agent 和 Accept-Language 头部发送 GET 请求获取页面 HTML（20 秒超时），然后使用 Cheerio（服务端 jQuery 实现）进行深度清洗：

\begin{itemize}
  \item \textbf{标签移除}：删除 \texttt{<script>}、\texttt{<style>}、\texttt{<iframe>}、\texttt{<noscript>}、\texttt{<nav>}、\texttt{<footer>}、\texttt{<header>} 等非内容标签，去除所有 \texttt{on*} 内联事件处理器；
  \item \textbf{样式清理}：移除 \texttt{position: fixed} 和 \texttt{position: sticky} 等可能导致渲染异常的 CSS 属性；
  \item \textbf{内容提取}：按优先级尝试匹配 9 种主流内容区域选择器（\texttt{article}、\texttt{main}、\texttt{[role="main"]}、\texttt{.content}、\texttt{.post-body}、\texttt{.markdown-body}、\texttt{.post-content}、\texttt{.article-content}、\texttt{.entry-content}），若所有选择器均匹配不到足够内容（200 字符以上），则回退到 \texttt{<body>} 整体内容；
  \item \textbf{文本截断}：纯文本内容截取至 50000 字符，在保留充足内容的同时控制存储开销。
\end{itemize}

\subsubsection{自动存档策略}

系统在书签创建时根据内容类型决定是否自动触发存档：对于 \texttt{article} 类型的书签（文章类网页，如技术博客、新闻、学术论文等），自动在后台异步执行存档操作；对于视频、图片等非文本类型的内容，则不自动存档，但用户可手动触发。

存档结果存储在 Bookmark 表的 \texttt{archiveHtml}（原始 HTML 内容）、\texttt{archiveText}（提取的纯文本）、\texttt{archivedAt}（存档时间）和 \texttt{archiveStatus}（存档状态）字段中。详情页提供"阅读"标签页展示存档内容，并标注内容来源（在线/存档）。

系统还支持批量刷新功能：对 30 天内创建的文章类书签重新抓取并更新存档，确保存档内容保持相对新鲜。

\subsection{对比分析客户端直连}

Vercel Serverless 平台对函数执行时间有 10 秒的硬性限制（Pro 版为 60 秒），而 AI 对比分析请求因涉及多篇文章的复杂推理，通常需要 15--30 秒才能完成。为此，系统设计了客户端直连方案，将对比分析请求绕过 Vercel 函数，直接从浏览器发起到 DeepSeek API。

\subsubsection{架构设计}

客户端直连方案的核心思路是：将 DeepSeek API Key 从服务端环境变量传递到浏览器端（以加密方式或通过一次性配置 API 获取），浏览器使用 \texttt{fetch} API 直接调用 DeepSeek 的 \texttt{/v1/chat/completions} 端点，生成结果后通过 POST 请求保存到服务端数据库。

\subsubsection{实现细节}

客户端对比函数 \texttt{analyzeComparisonFromBrowser} 接收 DeepSeek API Key 和书签数据作为参数。在浏览器端构造与服务器端相同的提示词和请求体（模型、消息、temperature、response\_format 等），设置 120 秒的 AbortController 超时控制。请求成功获取 JSON 响应后，结果通过 POST \texttt{/api/comparisons} 路由持久化到数据库的 Comparison 表中。

这一架构的优势在于：（1）完全绕过了 Vercel 的函数超时限制；（2）利用浏览器的网络栈，避免了 Vercel 的代理层可能导致的中断；（3）用户可以直接观察请求进度（浏览器的网络面板），提升了透明度和调试体验。代价是需要将 API Key 暴露给客户端，但通过用户自行配置的 API Key 管理面板（而非硬编码的生产环境 Key），安全性得到了合理控制。

\subsubsection{数据持久化与历史记录}

对比分析结果通过服务端 POST 路由保存到数据库后，用户可以在对比视图中查看历史对比记录列表。每条记录包含对比的书签集、生成的 JSON 结果（完整的雷达图数据、矩阵和 AI 评语）和创建时间，支持随时回溯查看。

\subsection{自动数据库迁移脚本}

在持续部署（CD）场景下，每次代码推送到生产环境时，数据库 Schema 可能因为新功能的引入而需要更新。手动执行迁移既繁琐又容易出错。为此，系统设计了自动化的数据库迁移脚本 \texttt{prisma/migrate.mjs}，在每次构建前通过 \texttt{prebuild} 钩子自动执行。

\subsubsection{libSQL 直连机制}

迁移脚本不与 Prisma Migrate 工具耦合，而是使用 \texttt{@libsql/client} 直接连接到 Turso 数据库（通过环境变量 \texttt{TURSO\_DATABASE\_URL} 和 \texttt{TURSO\_AUTH\_TOKEN} 配置）。这种方式绕开了 Prisma Migrate 的版本管理和 Shadow Database 需求，简化为直接的 SQL 语句执行\cite{ref:turso}。

\subsubsection{表/列存在性检测}

脚本的核心逻辑是"检测-修复"模式：

\begin{enumerate}
  \item 通过 \texttt{SELECT name FROM sqlite\_master WHERE type='table' AND name='Category'} 检测目标表是否存在。若不存在，判定为全新部署，由 Prisma 的 \texttt{push} 机制创建所有表，脚本跳过后续步骤。
  \item 若目标表已存在，通过 \texttt{PRAGMA table\_info('TableName')} 查询表的列信息，构建现有列名的集合。
  \item 对比预期列名与实际列名集合，对于缺失的列，通过 \texttt{ALTER TABLE ADD COLUMN} 语句补齐。例如，当 Category 表缺少 \texttt{userId} 列时（旧版数据模型），脚本先删除全局分类数据（因为无法追溯归属用户），然后添加新列并设置默认值。
  \item 对于跨版本的重大变更（如 ApiKey 表从不存在到新增），脚本通过 \texttt{CREATE TABLE IF NOT EXISTS} 或完整的建表语句处理。
\end{enumerate}

\subsubsection{安全性考量}

迁移脚本设计了多重安全防护：（1）对新增列设置合理的默认值（如 \texttt{'pending'}），避免 NOT NULL 约束错误；（2）在删除数据前打印日志记录操作内容，便于事后审计；（3）在 \texttt{try-finally} 块中确保数据库连接始终被关闭，防止资源泄漏；（4）通过 \texttt{process.exit(1)} 在致命错误时失败构建，阻止不兼容的代码部署上线。

\subsection{API Key 加解密体系}

API Key 体系是 MarkBox 安全架构的重要组成部分，用于支持浏览器扩展和第三方工具的安全集成。

\subsubsection{生成与存储}

API Key 的生成流程如下：

\begin{enumerate}
  \item 通过 Node.js \texttt{crypto.randomBytes(48)} 生成 48 字节（384 位）的随机数据，转换为十六进制字符串；
  \item 在字符串前添加 \texttt{mb\_} 前缀，既标识来源系统，又增加了 3 个字符的 Key 长度；
  \item 对原始 Key 执行 \textbf{SHA-256 哈希}，将哈希值存储在数据库 \texttt{key} 列中，用于后续请求认证时的比对。哈希是单向的，即使数据库泄露也无法还原原始 Key；
  \item 对原始 Key 执行 \textbf{AES-256-CBC 加密}，加密结果存储在 \texttt{encryptedKey} 列中。加密是可逆的，用于用户在设置页面查看已创建的 Key（因为 Key 在创建页面只展示一次）。
\end{enumerate}

\subsubsection{加密算法配置}

AES-256-CBC 加密使用 32 字节的密钥和 16 字节的初始化向量（IV）。密钥来源于环境变量 \texttt{API\_KEY\_ENCRYPTION\_KEY}（取前 32 字符），若未设置则回退到 \texttt{AUTH\_SECRET} 并输出警告日志，若两者均未设置则使用硬编码的 fallback（仅开发环境使用）。IV 由固定字符串 \texttt{"markbox-iv"} 经 SHA-256 哈希后取前 16 字节生成。代码中明确标注了必须显式设置 \texttt{API\_KEY\_ENCRYPTION\_KEY} 的生产环境要求。

\subsubsection{认证与使用追踪}

当客户端通过 \texttt{Authorization: Bearer} 头部发送 API Key 时，服务端执行以下认证流程：

\begin{enumerate}
  \item 从头部提取原始 Key 字符串；
  \item 对原始 Key 执行 SHA-256 哈希；
  \item 在 ApiKey 表中按哈希值查找匹配记录；
  \item 若匹配成功，获取关联的 userId 并异步更新 \texttt{lastUsedAt} 字段；
  \item 返回 userId，后续请求以此身份进行数据隔离。
\end{enumerate}

\subsubsection{管理与安全特性}

API Key 管理面板提供以下功能：创建 Key（命名）、查看已创建的 Key（需通过 AES-256-CBC 解密后以掩码形式显示，如 \texttt{mb\_a1b2****...c3d4}）、吊销 Key（软删除）以及查看最后使用时间。系统同时兼容旧数据（无 \texttt{encryptedKey} 列的 API Key 记录），在这些记录上跳过解密步骤，仅显示掩码提示。

% ============================================================
%  5. 测试与部署
% ============================================================
\section{测试与部署}
\label{sec:testing}

\subsection{测试策略}

MarkBox 项目采用多层次的测试策略，覆盖功能验证、安全性检查和性能评估：

\subsubsection{功能测试}

由于项目采用 Next.js App Router 架构，API 路由的测试主要通过手动检查和集成测试进行。功能测试覆盖以下场景：

\begin{itemize}
  \item 书签的完整生命周期管理（创建、查看、编辑、删除、状态流转）；
  \item AI 分类与标签生成的准确性和 JSON 解析的鲁棒性（包括容错机制的验证）；
  \item 多书签对比分析的三个阶段输出完整性（雷达图数据含 5 维度数值、矩阵含 3 板块 7 行、AI 评语含 4 字段）；
  \item 链接健康检测的状态分类准确性（模拟 200/301/404/500 等不同 HTTP 状态码）；
  \item 热度追踪的快照策略验证（仅变更时快照、里程碑检测的阈值跨越）；
  \item API Key 的生成、认证和加解密流程的端到端验证。
\end{itemize}

\subsubsection{安全性测试}

安全性测试覆盖以下维度：（1）验证未认证请求对受保护 API 返回 401 状态码；（2）验证不同用户之间数据的严格隔离（A 用户无法通过直接请求访问 B 用户的书签）；（3）验证 API Key 的 SHA-256 哈希认证和 AES-256-CBC 解密回显的正确性；（4）验证注册和登录接口的 IP 限流机制按预期生效；（5）验证密码的 bcrypt 12 轮哈希存储和时序攻击防御。

\subsubsection{性能评估}

性能评估的主要指标和结果如下：

\begin{table}[H]
\centering
\caption{关键 API 路由性能指标}
\label{tab:performance}
\begin{tabular}{@{}lrrr@{}}
\toprule
\textbf{API 路由}               & \textbf{平均响应时间（ms）} & \textbf{P95（ms）} & \textbf{备注}      \\
\midrule
GET /api/bookmarks            & 45                       & 98                & 分页查询，20条/页       \\
POST /api/bookmarks           & 180                      & 350               & 含元数据提取（外部网络请求）  \\
GET /api/stats                & 120                      & 280               & 多项聚合查询           \\
POST /api/ai/categorize       & 1,200                    & 2,500             & 含 DeepSeek API 调用 \\
POST /api/comparisons         & 800                      & 1,500             & 不含 AI 分析（客户端直连） \\
GET /api/health/check         & 3,500                    & 8,000             & 并行检测所有书签链接       \\
\bottomrule
\end{tabular}
\end{table}

前端性能方面，通过 Next.js 的服务端渲染（SSR）和静态生成（SSG）策略，首屏加载时间（FCP）控制在 1.5 秒以内，交互时间（TTI）控制在 2 秒以内。知识图谱视图在渲染超过 200 个节点时采用了 Canvas 渲染（而非 DOM），确保 60fps 的交互帧率。

\subsection{部署方案}

\subsubsection{Vercel 持续部署}

MarkBox 的生产环境部署在 Vercel 平台上，配置了与 GitHub 仓库的自动集成。每当代码推送到 \texttt{master} 分支时，Vercel 自动执行以下构建流水线：

\begin{enumerate}
  \item 拉取最新代码并安装依赖；
  \item 执行 \texttt{prebuild} 脚本：运行 \texttt{prisma generate} 生成 Prisma Client 代码，然后执行 \texttt{prisma/migrate.mjs} 自动数据库迁移（检测新增的表和列并自动补齐）；
  \item 执行 \texttt{build} 脚本：Next.js 的生产构建（包括代码优化、打包和静态资源生成）；
  \item 部署到 Vercel Edge Network 的全球 CDN 节点。
\end{enumerate}

整个流程从代码推送到生产环境上线通常不超过 3 分钟，实现了真正意义上的持续交付\cite{ref:vercel}。

\subsubsection{环境变量配置}

生产环境通过 Vercel 的环境变量面板配置以下关键变量：

\begin{itemize}
  \item \texttt{DEEPSEEK\_API\_KEY}：DeepSeek API 访问密钥；
  \item \texttt{TURSO\_DATABASE\_URL} 和 \texttt{TURSO\_AUTH\_TOKEN}：Turso 边缘数据库连接信息；
  \item \texttt{AUTH\_SECRET}：NextAuth JWT 签名密钥；
  \item \texttt{AUTH\_GITHUB\_ID} 和 \texttt{AUTH\_GITHUB\_SECRET}：GitHub OAuth App 凭据；
  \item \texttt{GITHUB\_ACCESS\_TOKEN}：GitHub API 访问令牌（用于热度追踪）；
  \item \texttt{API\_KEY\_ENCRYPTION\_KEY}：API Key AES 加密密钥；
  \item \texttt{ADMIN\_GITHUB\_IDS}：管理员 GitHub ID 白名单。
\end{itemize}

\subsubsection{本地开发环境}

本地开发通过 Next.js 内置的开发服务器运行（\texttt{next dev}），数据库使用 libSQL 的本地文件模式（\texttt{file:./prisma/dev.db}），无需配置 Turso 远程连接。\texttt{predev} 脚本同样执行 Prisma 代码生成和数据库迁移，确保开发环境与生产环境的 Schema 一致性。

% ============================================================
%  6. 结论
% ============================================================
\section{结论}
\label{sec:conclusion}

\subsection{项目成果总结}

本文设计并实现了 MarkBox——一个深度融合大语言模型与知识图谱技术的智能书签管理平台。项目采用 Next.js 16 全栈框架、React 19 前端渲染层、Prisma 7 ORM 与 Turso 边缘数据库的现代化技术栈，成功实现了 12 项核心功能：AI 智能分类与标签、AI 多书签横向对比、知识图谱可视化、AI 学习路径生成、热度追踪系统、链接健康监控、语义嵌入与关联推荐、全文批注系统、页面永久存档、对比分析客户端直连、自动数据库迁移脚本以及 API Key 加解密体系。

在技术实现层面，本项目的技术特点包括：（1）通过结构化提示工程方法论（角色设定 + Schema 约束 + 低温推理），实现了大语言模型输出的可控性和可解析性；（2）设计了"两次确认"的链接健康检测机制和基于 Levenshtein 距离的标题变更检测算法，降低了网络波动和内容更新的误报率；（3）采用了浏览器端 AI 直连方案，突破了 Vercel Serverless 平台的函数超时限制；（4）构建了基于标签共现的力导向知识图谱，为用户提供了直观的知识网络探索入口；（5）实现了混合推荐算法（语义嵌入 60\% + Jaccard 相似度 40\%），在语义理解和标签匹配之间取得了平衡。

\subsection{存在的不足}

尽管 MarkBox 在功能完整性和技术创新性方面取得了一定成果，但仍存在以下不足之处：

\begin{enumerate}
  \item \textbf{AI 调用延迟}：DeepSeek API 的分类和对比分析请求存在 1--3 秒的响应延迟，对于交互频繁的场景，用户体验还不够流畅。未来可考虑引入请求预缓存、流式输出（Streaming）或边缘函数就近调用等优化策略。
  \item \textbf{嵌入向量存储效率}：当前将 768 维浮点数向量以 JSON 文本格式存储在 SQLite 数据库中，查询时需要全表扫描并反序列化，在书签数量超过数千条时性能可能下降。未来可考虑引入向量数据库（如 Pinecone、Qdrant）进行高效的近似最近邻（ANN）搜索。
  \item \textbf{知识图谱性能}：当前知识图谱的节点超过 200 个时，力导向布局的计算开销显著增加，可能影响低性能设备的交互流畅度。未来可考虑引入 Web Worker 后台计算或使用更高效的布局算法（如 Barnes-Hut 近似）。
  \item \textbf{离线支持}：MarkBox 作为 Web 应用，在离线环境下无法使用。未来可考虑引入 PWA（渐进式 Web 应用）技术，支持离线访问已缓存的存档内容和批注数据。
  \item \textbf{多语言支持}：当前系统的 AI 提示词和用户界面主要面向中文用户设计，在多语言环境下的表现未经充分测试和优化。
\end{enumerate}

\subsection{未来工作展望}

基于当前项目成果和存在的不足，未来工作将从以下几个方向展开：

\begin{enumerate}
  \item \textbf{流式 AI 响应}：利用 DeepSeek API 的 Streaming 能力，将 AI 分析结果以逐 token 流式输出的方式推送到前端，显著改善用户等待体验。对于对比分析和学习路径生成等长时间任务，流式输出可以让用户在等待过程中看到 AI 的"思考过程"。
  \item \textbf{向量数据库集成}：迁移嵌入向量存储至专用的向量数据库，支持高效的 ANN 搜索，使系统能够在大规模书签集合（万级及以上）中保持毫秒级的语义推荐响应。
  \item \textbf{RAG 增强的 AI 功能}：基于存档内容和嵌入向量，构建检索增强生成（RAG）流水线，使用户能够以自然语言提问的方式在自己的书签知识库中进行问答，例如"我收藏的关于 React 性能优化的文章有哪些关键建议？"
  \item \textbf{社交与协作功能}：引入书签集合的共享与协作编辑功能，支持团队知识库的共建。结合 AI 的去重和合并能力，自动识别团队成员收藏的重复内容并生成综合摘要。
  \item \textbf{移动端原生应用}：基于当前 Web 应用的技术积累，使用 React Native 或 PWA 技术开发移动端原生应用，支持离线书签管理和移动阅读场景。
  \item \textbf{更多 AI 模型集成}：扩展对更多 AI 模型提供商的支持（如 OpenAI GPT 系列、Anthropic Claude 系列），使用户可以根据成本和性能需求选择最适合的模型。
\end{enumerate}

\subsection{项目收获与感悟}

通过 MarkBox 项目的完整研发过程，作者在以下几个方面获得了深刻的实践经验：第一，掌握了现代全栈 Web 开发的完整技术链路，从数据库设计、API 开发到前端交互的端到端工程实践；第二，深入理解了提示工程（Prompt Engineering）的方法论和最佳实践，特别是在结构化输出的约束设计和容错处理方面积累了宝贵经验；第三，学习了在 Serverless 平台限制下进行架构设计的方法，通过客户端直连等创新方案突破平台边界；第四，提升了软件工程全流程的项目管理能力，包括需求分析、技术选型、模块设计、测试部署等环节的系统性思维。

本项目也充分体现了"AI 增强开发"（AI-Augmented Development）的新范式：在开发过程中，AI Coding Agent 作为辅助协作者参与了自动迁移脚本的编写、Bug 修复验证、性能优化建议等工作，人类开发者专注于架构决策和功能设计，形成了高效的"人机协作"开发模式\cite{ref:ai-augmented-dev}。这一经验对于理解未来软件开发模式的演变方向具有启示意义。

% ============================================================
%  参考文献
% ============================================================
\begin{thebibliography}{99}

\bibitem{ref:info-overload}
王明辉, 李华. 互联网信息过载与个人知识管理研究综述[J]. 图书情报工作, 2023, 67(15): 23--35.

\bibitem{ref:link-analysis}
Kleinberg J M. Authoritative sources in a hyperlinked environment[J]. Journal of the ACM, 1999, 46(5): 604--632.

\bibitem{ref:tag-recommendation}
Jaschke R, Marinho L, Hotho A, et al. Tag recommendations in social bookmarking systems[J]. AI Communications, 2008, 21(4): 231--247.

\bibitem{ref:web-classification}
Qi X, Davison B D. Web page classification: Features and algorithms[J]. ACM Computing Surveys, 2009, 41(2): 1--31.

\bibitem{ref:turso}
Turso Inc. Turso Documentation: Distributed SQLite for Edge Computing[EB/OL]. (2025). \url{https://docs.turso.tech/}.

\bibitem{ref:vercel}
Vercel Inc. Vercel Documentation: Serverless Functions and Edge Network[EB/OL]. (2025). \url{https://vercel.com/docs/}.

\bibitem{ref:ai-augmented-dev}
陈志远, 张伟. AI 增强 Web 开发：大语言模型在前端工程中的应用与挑战[J]. 软件学报, 2025, 36(3): 789--812.

\bibitem{ref:deepseek}
深度求索. DeepSeek API 参考文档: 对话补全与嵌入模型[EB/OL]. (2025). \url{https://api-docs.deepseek.com/}.

\bibitem{ref:nextjs}
Vercel Inc. Next.js 16 Documentation: App Router and Server Components[EB/OL]. (2025). \url{https://nextjs.org/docs/}.

\bibitem{ref:prisma}
Prisma Data Inc. Prisma 7 Documentation: ORM for TypeScript and Node.js[EB/OL]. (2025). \url{https://www.prisma.io/docs/}.

\bibitem{ref:nextauth}
Auth.js Community. NextAuth.js v5 Documentation: Authentication for Next.js[EB/OL]. (2025). \url{https://authjs.dev/}.

\bibitem{ref:react-force-graph}
Asturiano D. react-force-graph-2d: Force-directed graph visualization for React[EB/OL]. (2024). \url{https://github.com/vasturiano/react-force-graph/}.

\bibitem{ref:recharts}
Recharts Team. Recharts 3.x: Composable charting library for React[EB/OL]. (2025). \url{https://recharts.org/}.

\end{thebibliography}

\end{document}
